% !TeX root = ../navier-stokes-manuscript.tex
% ============================================================
% 1.2 Historical Background
% ============================================================
\subsection{Historical Background}\label{sub:HistoricalEstimates}

A smooth local solution gives a genuine Navier--Stokes fluid for a short time, while the Clay problem asks whether that same solution can cross every finite time.
The history of the problem can be read as a sequence of attempts to close that gap.
Each major estimate controls one feature of the motion, reveals what remains free to concentrate, and leaves a sharper version of the continuation question behind it.
The order below follows that tightening of the problem, with publication dates serving as secondary orientation.

\divider{Leray's Global Energy Control}

The first decisive calculation asks what the equations preserve before any derivative is placed under pointwise control.
For a smooth solution, take the spatial inner product of the momentum equation with \(u\).
Incompressibility turns the transport term into a spatial divergence and makes the pressure term orthogonal to the velocity, while integration by parts turns viscosity into dissipation.
The result is
\[
  \lVert u(t)\rVert_{L^2}^{2}
  +2\nu\int_0^t\lVert\nabla u(s)\rVert_{L^2}^{2}\,ds
  =
  \lVert u_0\rVert_{L^2}^{2}.
\]
The velocity stays bounded in \(L^2\) at each time, and viscosity pays for the total \(L^2\) gradient cost accumulated over the interval.

Leray used regularized equations, compactness, and weak lower semicontinuity to carry the corresponding inequality into a global weak solution \parencite{Leray1934}.
On every finite interval, a Leray solution satisfies
\[
  u\in L^\infty(0,T;L^2(\Omega))
  \cap
  L^2(0,T;H^1(\Omega)).
\]
This is genuine global existence in the energy class.
It does not settle the Clay problem because smoothness and uniqueness are not known for arbitrary three-dimensional Leray solutions.

The two terms in the estimate also carry different information in time.
The \(L^2\) size of \(u\) is bounded at every instant, whereas the \(L^2\) size of \(\nabla u\) is controlled only after time integration.
A function can have finite area under its graph and still form arbitrarily tall peaks on arbitrarily short intervals.
The energy inequality leaves open the possibility that derivative activity concentrates near a finite time even though the total kinetic energy remains finite.

\divider{Enstrophy and Vortex Stretching}

For a divergence-free field on \(\mathbb R^3\) with sufficient decay, or on \(\mathbb T^3\),
\[
  \lVert\nabla u\rVert_{L^2}^{2}
  =
  \lVert\omega\rVert_{L^2}^{2},
  \qquad
  \omega:=\nabla\times u.
\]
The quantity \(\lVert\omega\rVert_{L^2}^{2}\) is the enstrophy, so Leray's estimate controls the area under the enstrophy curve without controlling its height.

The vorticity equation shows why this distinction becomes serious in three dimensions:
\[
  \partial_t\omega+(u\cdot\nabla)\omega
  =
  (\omega\cdot\nabla)u+\nu\Delta\omega.
\]
Transport moves vorticity with the fluid, and viscosity diffuses sharp spatial variation.
The remaining term, \((\omega\cdot\nabla)u\), stretches and tilts vorticity through the surrounding velocity gradient.
Its contribution to the enstrophy balance is
\[
  \frac12\frac{d}{dt}\lVert\omega\rVert_{L^2}^{2}
  +\nu\lVert\nabla\omega\rVert_{L^2}^{2}
  =
  \int_\Omega (\omega\cdot\nabla)u\cdot\omega\,dx.
\]
The stretching integral has no fixed sign.
Viscosity continues to oppose spatial concentration, while the energy estimate alone gives no pointwise-in-time comparison proving that diffusion always defeats stretching.

This is the first clear shape of the three-dimensional obstruction.
The velocity cannot escape the global energy bound, yet the derivatives that determine classical smoothness may rise through narrow peaks that the energy balance does not see individually.
This leaves a precise question: how much derivative control would be enough to continue the smooth solution?

\divider{Sobolev Regularity and Continuation}

Sobolev spaces answer that question by measuring velocity and spatial derivatives in one norm.
For an integer \(m\geq0\),
\[
  H^m(\Omega)
  =
  \left\{
    v:
    \partial^\alpha v\in L^2(\Omega)
    \text{ for every }|\alpha|\leq m
  \right\},
\]
with the usual weak interpretation of the derivatives and the Fourier definition for fractional orders.
An \(L^p\) norm records the spatial size and concentration of the function itself, while an \(H^m\) norm also records its variation through order \(m\).

Sobolev embedding converts sufficiently high \(L^2\)-derivative control into ordinary continuous derivatives.
In three dimensions \parencite[Theorem~7.28]{Cerda2010},
\[
  H^s(\Omega)\hookrightarrow C^k(\Omega)
  \qquad
  \text{whenever}
  \qquad
  s>k+\frac32.
\]
One \(H^s\) bound buys only finitely many continuous derivatives.
For example, \(s>5/2\) gives
\[
  H^s(\Omega)\hookrightarrow W^{1,\infty}(\Omega),
\]
so the velocity and its first spatial derivative are bounded and continuous.

The threshold \(s>5/2\) is useful because the Lipschitz norm of \(u\) controls the way transport amplifies every higher derivative already present in the smooth datum.
For each fixed finite \(m\geq s\), the standard differentiated estimate has the form
\[
  \frac12\frac{d}{dt}\lVert u\rVert_{H^m}^{2}
  +\nu\lVert\nabla u\rVert_{H^m}^{2}
  \leq
  C_m\lVert\nabla u\rVert_{L^\infty}\lVert u\rVert_{H^m}^{2}.
\]
Once \(\lVert\nabla u\rVert_{L^\infty}\) is integrable in time, Gronwall's inequality propagates every fixed finite Sobolev order from the smooth initial datum.
The constant may depend on \(m\), which is enough because smoothness asks for every finite order separately and imposes no one numerical bound across the infinite family.
Applying Sobolev embedding at successively higher orders then recovers every continuous spatial derivative.

The pressure follows from the velocity through the elliptic relation
\[
  -\Delta p
  =
  \sum_{i,j=1}^{3}\partial_i\partial_j(u_i u_j),
\]
with zero mean on the torus or the usual decay normalization on \(\mathbb R^3\).
Once all spatial derivatives of \(u\) and \(p\) are controlled, the momentum equation
\[
  \partial_tu
  =
  \nu\Delta u-(u\cdot\nabla)u-\nabla p
\]
recovers the first time derivative.
Repeatedly differentiating the pressure relation and the momentum equation then recovers every finite mixed spatial and temporal derivative.
This is the continuation argument from a strong Sobolev bound to the \(C^\infty\) velocity--pressure pair required by the problem.

The local theories of Fujita--Kato and Kato supply the restart mechanism behind this argument \parencite{FujitaKato1964,Kato1984}.
Suppose a smooth solution exists on \([0,T)\) and, for one \(s>5/2\),
\[
  \sup_{0\leq t<T}\lVert u(t)\rVert_{H^s}\leq M.
\]
Choose an actual time \(t_0<T\) close to \(T\).
Local theory restarted from \(u(t_0)\) has a lifespan bounded below in terms of \(s\), \(M\), and \(\nu\), so a sufficiently late choice of \(t_0\) produces a solution beyond \(T\).
Uniqueness identifies the restarted solution with the original one on their overlap.
Thus a finite maximal smooth time \(T_*\) would force every continuation-grade \(H^s\) norm to become unbounded as \(t\uparrow T_*\).

This criterion identifies the exact kind of estimate that would close the problem.
The energy inequality supplies
\[
  u\in L^\infty_tL^2_x\cap L^2_tH^1_x,
\]
while continuation asks for pointwise-in-time control at a substantially stronger derivative level.
The distance between those two statements is the analytic room in which a finite-time singularity could still hide.

\divider{Partial Regularity and the Shape of a Singular Set}

Continuation theory describes what must diverge along a maximal classical solution, while the global weak theory asks where regularity can fail after that classical description is no longer available.
Caffarelli, Kohn, and Nirenberg studied suitable weak solutions through the local energy inequality on shrinking parabolic cylinders \parencite{CaffarelliKohnNirenberg1982}.
For \(z_0=(x_0,t_0)\), write
\[
  Q_r(z_0)
  :=
  B_r(x_0)\times(t_0-r^2,t_0).
\]
Their epsilon-regularity mechanism says that sufficiently small scale-invariant velocity and pressure quantities on \(Q_r(z_0)\) force regularity on a smaller cylinder.

At a singular point, the contrapositive applies at every sufficiently small scale.
Every shrinking cylinder centred there must retain a definite amount of scale-invariant concentration, because one cylinder below the threshold would already make the point regular.
A covering argument then confines the possible singular set to a set of zero one-dimensional parabolic Hausdorff measure.
This is an exceptionally small set, yet even one singular point would defeat global smoothness.

Partial regularity changes the picture of failure.
A counterexample cannot be a diffuse loss of control spread harmlessly across the fluid.
It must concentrate velocity and pressure through every smaller parabolic neighbourhood of at least one spacetime point.
That conclusion creates the need for criteria that measure spatial intensity and temporal persistence together.

\divider{Scale-Critical Continuation Criteria}

The Ladyzhenskaya--Prodi--Serrin criteria provide that joint measurement \parencite{Prodi1959,Serrin1962,Serrin1963,Ladyzhenskaya1967}.
They give regularity through \(T\) whenever
\[
  u\in L^p(0,T;L^q(\Omega)),
  \qquad
  \frac{2}{p}+\frac{3}{q}\leq1,
  \qquad
  q>3.
\]
The exponent \(q\) measures the spatial intensity of a concentration, and \(p\) measures how long that intensity can persist.

At one derivative, interpolation and Young's inequality lead to
\[
  \frac{d}{dt}\lVert\nabla u\rVert_{L^2}^{2}
  +\nu\lVert\Delta u\rVert_{L^2}^{2}
  \leq
  C_{\nu,q}
  \lVert u\rVert_{L^q}^{\frac{2q}{q-3}}
  \lVert\nabla u\rVert_{L^2}^{2}.
\]
The mixed-norm condition makes the coefficient on the right integrable in time, so Gronwall's inequality keeps the enstrophy finite and continuation crosses \(T\).
At the equality
\[
  \frac{2}{p}+\frac{3}{q}=1,
\]
the criterion is invariant under the Navier--Stokes scaling
\[
  u_\lambda(x,t)
  :=
  \lambda u(\lambda x,\lambda^2t).
\]
These critical exponents match exactly the balance between stronger spatial concentration and shorter parabolic duration.

The direct estimate above covers every \(q>3\) and degenerates as \(q\downarrow3\).
The endpoint \(q=3\) is especially important because
\[
  \lVert u_\lambda(t)\rVert_{L^3}
  =
  \lVert u(\lambda^2t)\rVert_{L^3}.
\]
Magnifying a suspected singularity to unit scale does not weaken the norm used to measure it.

Escauriaza, Seregin, and \v{S}ver{\'a}k proved that boundedness in \(L^\infty_tL^3_x\) still excludes a finite singular time \parencite{EscauriazaSereginSverak2003}.
Their proof assumes a first singular time, rescales around the concentration, extracts a nontrivial ancient limiting solution, and uses backward uniqueness and unique continuation to force that limit to vanish.
The contradiction shows that
\[
  u\in L^\infty(0,T;L^3(\mathbb R^3))
\]
is sufficient for regularity through \(T\).

Seregin subsequently showed, for the smooth compactly supported data in his whole-space result, that a finite maximal time forces full divergence of the critical norm \parencite{Seregin2012}:
\[
  \lim_{t\uparrow T_*}\lVert u(t)\rVert_{L^3(\mathbb R^3)}
  =
  \infty.
\]
The critical norm cannot spike along a thin sequence of times and repeatedly return below a fixed level.
Every finite level must eventually be left behind as the alleged endpoint approaches.

Tao later quantified part of this conclusion \parencite{Tao2021Quantitative}.
Along infinitely many times approaching a finite maximal time, the \(L^3\) norm must exceed a positive power of a triple logarithm of the inverse time remaining.
This rate is weak compared with a power law, yet it proves that a finite-time singularity would require a genuine cascade through successively smaller scales.

\divider{Where the Problem Sits Today}

The established theory entering this paper replaces an undefined breakdown with a sharply constrained alternative.
If \(T_*<\infty\), the kinetic energy remains bounded, while continuation-grade Sobolev norms and the critical \(L^3\) norm must diverge.
Any singular point must retain scale-invariant concentration in every sufficiently small parabolic cylinder around it.
The concentration may accelerate through shorter and shorter time intervals, because parabolic scale reduction also shortens the available time scale.

The basic scaling already explains why the global energy estimate cannot rule this out.
At corresponding times,
\[
  \lVert u_\lambda(t)\rVert_{L^3}
  =
  \lVert u(\lambda^2t)\rVert_{L^3},
  \qquad
  \lVert u_\lambda(t)\rVert_{L^2}^{2}
  =
  \lambda^{-1}\lVert u(\lambda^2t)\rVert_{L^2}^{2}.
\]
Compression to a shorter length scale preserves the critical \(L^3\) size while reducing the \(L^2\) energy seen at that scale.
A scale-independent ceiling on the critical norm cannot be read directly from the global energy bound.

A simple frequency-doubling picture makes the remaining timing problem visible.
If an active frequency doubles at each stage, its spatial scale halves and its parabolic time scale quarters.
The resulting durations form a convergent geometric series, so positive duration at every stage still leaves room for infinitely many stages before one finite time.
This is the Zeno-cascade picture behind the endpoint threat.
It is a heuristic picture of the allowed timing, not a construction of a Navier--Stokes singularity.

The public frontier has isolated the exact burden.
Known theory says what must diverge, how concentrated the divergence must become, and which bounds would exclude it.
It does not derive from arbitrary smooth data, fixed viscosity, and the equations a continuation-grade bound that remains finite all the way to every candidate endpoint.
That is the direct estimate a Gold-standard proof would have to supply.
